\documentclass[11pt]{article}

\usepackage[T1]{fontenc}
\usepackage[margin=2.5cm]{geometry}
\usepackage{amsmath,amssymb,amsthm}
\usepackage{booktabs,array}
\usepackage{hyperref}

\title{ML behaviour of the univariateML estimators:\\ Hellinger closure, escapes, and what the fitters do}
\author{}
\date{\today}

\newcommand{\E}{\mathbb{E}}
\newcommand{\tr}{\operatorname{tr}}
\newcommand{\Var}{\operatorname{Var}}
\newcommand{\Cov}{\operatorname{Cov}}
\newcommand{\vech}{\operatorname{vech}}
\newcommand{\argmin}{\operatorname*{arg\,min}}
\newcommand{\argmax}{\operatorname*{arg\,max}}
\newcommand{\plim}{\operatorname*{plim}}
\newcommand{\cF}{\mathcal{F}}
\newcommand{\code}[1]{\texttt{#1}}

\theoremstyle{plain}
\newtheorem{theorem}{Theorem}
\newtheorem{lemma}{Lemma}
\newtheorem{corollary}{Corollary}
\newtheorem{proposition}{Proposition}
\theoremstyle{definition}
\newtheorem{definition}{Definition}
\theoremstyle{remark}
\newtheorem*{remark}{Remark}

\begin{document}
\maketitle

\section*{Purpose}

Which of the 42 families are not Hellinger-closed, where the maximum
likelihood estimator (MLE) escapes to a limit model outside the family, and
what the current fitters do in those cases. First pass (October 2026),
updated the same day with a literature and derivation report
(\code{dev/research/ml-behaviour-escape-theory.md}) and a sweep of all 42 fitters
against a slow oracle (\code{dev/research/ml-behaviour-sweep.\{R,md\}}); the first
probe is \code{dev/research/ml-behaviour-probe.R}. Each row says how firm it is. The
per-distribution conclusions are meant to end up in each distribution's own
file.

Main findings.
\begin{enumerate}
  \item Six families escape with a ``use \code{mlX}'' message: lomax,
    gompertz, nbinom, binom, burr, and (since 2026-10-28) invburr.
  \item std and snorm escaped silently; since 2026-10-28 they report it.
    std stops with ``use \code{mlnorm}'' (or an infinite-variance message
    for the $\nu\to2^+$ face); snorm is fitted exactly through a
    one-dimensional profile (Section~\ref{sec:snorm}) and reports the
    translated half-normal faces. sstd now compares its limiting faces and
    stops when a limit wins; its global comparisons remain numerical.
  \item ged and sged have \emph{no} MLE for any nonconstant sample: the
    likelihood is unbounded as the shape tends to $0$ \cite{tumlinson2015}.
    Their large-shape ``uniform'' fits are local maxima only.
  \item The sweep found invalid fits on ordinary data (zip, cauchy, logis,
    binom); see Section~\ref{sec:bugs}.
  \item lomax's old escape test $\overline{x^2} < 2\bar x^2$ (and
    invpareto's, on $1/x$) was only local; it rejected samples such as
    $(1, 100)$ that have an MLE. Since 2026-10-29 the fitters use the exact
    criterion, Theorem~\ref{thm:lomax} (Section~\ref{sec:lomax}).
\end{enumerate}

\section{Definitions}\label{sec:closure-frame}

For probability laws $P,Q$, write $H^2(P,Q)=\tfrac12\int(\sqrt p-\sqrt q)^2$ with respect to a common dominating measure. Hellinger-closed means that every Hellinger limit is a member. Earlier tables use ``closed apart from point masses'' as a separate convention, not as the definition of Hellinger closure.

Fix a sample $x=(x_1,\ldots,x_n)$ and a parameter space $\Theta$, including any holds and fit restrictions. Write $L_x(\theta)=\sum_i\log f_\theta(x_i)$, with R7a's pointwise endpoint convention and extended-real values. Choose a Hausdorff topology on a space $K$ containing $\Theta$ densely. The \emph{joint sample likelihood envelope} is
\[
 L_x^*(z)=\limsup_{\theta\to z,\,\theta\in\Theta} L_x(\theta).
\]
This is the upper semicontinuous closure of the joint density values at the data (equivalently the log of the limsup of their product). A separate limsup for each observation cannot in general be multiplied: the paths attaining them may be incompatible. The envelope may distinguish paths leading to the same law; density values at isolated points are not determined by a probability law. An escape occurs when the original supremum is not attained and a maximizing sequence approaches a closure point outside $\Theta$. If all maximizing sequences have this property, the escape is unavoidable. A named limit law is justified only when that path also converges to it, in Hellinger or weak topology as appropriate. Collapse may have infinite likelihood and a weak point-mass limit.

\begin{lemma}[Likelihood closure]\label{lem:closure}
Suppose $K$ is compact and metrizable, $\Theta$ is dense, and $L_x$ is upper semicontinuous on $\Theta$. Let $L_x^*$ be its upper semicontinuous envelope on $K$. Then
\[
 \sup_{\theta\in\Theta} L_x(\theta)=\max_{z\in K} L_x^*(z).
\]
If the maximum is finite and attained only outside $\Theta$, every maximizing sequence has a subsequence converging to a maximizing boundary point. The same conclusion holds when compactness is replaced by compactness of a set eventually containing some maximizing sequence; to apply it to \emph{every} maximizing sequence require this property for each one.
\end{lemma}
\begin{proof}
The envelope is upper semicontinuous and dominates $L_x$; upper semicontinuity on $\Theta$ gives equality there. Every envelope value is at most $S=\sup_\Theta L_x$, while density of $\Theta$ gives the reverse inequality for the supremum. An upper semicontinuous extended-real function attains its maximum on compact $K$. Compactness yields a convergent subsequence of a maximizing sequence, and upper semicontinuity makes its limit a maximizer. The local compactness variant follows by the same subsequence argument.
\end{proof}
This lemma does \emph{not} by itself identify $L_x^*(z)$ with the ordinary likelihood of a limiting law. That requires an additional, sample-specific compatibility condition: the joint envelope at each relevant boundary point equals the stated limit-model likelihood, including at endpoints, and the parametrization covers the candidate face. When the maximizing boundary lies on such a face and its likelihood maximum is attained, the family supremum equals that face's maximized log likelihood. A limiting law or upper semicontinuity alone does not prove that equality. Nor does weak convergence ensure tightness of arbitrary sequences. We assume maximizing sequences are tight as laws when invoking weak limit laws, and separately assume the parameter/likelihood compactness needed above; tightness alone neither supplies it nor excludes pointwise density spikes. No defective sub-probability limit was observed in the 98-family study; finite mixtures, where loss of mass can occur, are outside scope.

\section{Statistical meaning}\label{sec:statistical-meaning}
For the normalization above, $H^2\leq\mathrm{TV}\leq\sqrt2H$: these metrics generate the same topology. Continuous laws can converge weakly to a point mass while their TV distance from it remains one; discrete point masses can be Hellinger limits or attained members (R7). For proofs and examples see \href{../../univariateML-cpp/notes/limits.md}{the limits article}. Under misspecification an expected-log-likelihood maximizer, when well-defined and attained, is a KL projection onto the closure. Wald-type consistency of fitted laws for that projection additionally needs identifiability or a set-valued target, uniform convergence and appropriate compactness/tightness and integrability; Hellinger convergence alone gives neither KL convergence nor empirical likelihood convergence. A finite-sample escape does not establish those assumptions.

\paragraph{Family-level verification.} The classification tables record path arguments, code and probes, not a verification of the hypotheses of Lemma~\ref{lem:closure}. For all 98 families and their held subsets the joint-envelope continuity, compactness and tightness conditions are \emph{assumed}; individual paths carry the arguments and qualifications cited in their dossiers. The closed-form inverse-Gaussian mean comparison verifies non-reachability of its fixed-shape L\'evy face on finite positive samples, but not the global hypotheses of the lemma.

\section{Classification}

Abbreviations in the ``firm'' column: \emph{code} = condition implemented
and tested in the fitter; \emph{arg} = short limit argument in this note;
\emph{probe} = observed in a probe fit; \emph{conj} = conjecture.

\subsection{Escape, already caught with a pointer to the limit model}

\begin{center}
\small
\begin{tabular}{@{}l p{4.0cm} p{2.6cm} p{4.4cm} l@{}}
\toprule
id & escape & limit model & detection today & firm \\
\midrule
lomax & $\kappa\to\infty$, $\lambda\kappa$ fixed & exponential \cite{kk2003} &
  global profile search on $0<u<\bar u$ (Theorem~\ref{thm:lomax}), status
  \code{ExponentialLimit}; the old moment test was only local & code, arg \\
gompertz & $b \to 0$ & exponential & status \code{Overdispersed} & code, probe \\
nbinom & size $\to\infty$, mean fixed & Poisson &
  underdispersion check in the Rcpp adapter only & code \\
binom & size $\to\infty$, mean fixed & Poisson & status; message names \code{mlpois} & code \\
burr & shape1 $\to\infty$, scale$\cdot$shape1$^{-1/\text{shape2}}$ fixed & Weibull &
  face comparison, status \code{WeibullBoundary} & code, probe \\
burr & shape2 $\to\infty$, shape1$\cdot$shape2 fixed & Pareto, minimum = scale &
  face comparison; with the scale held at $s$: all $x_i \ge s$ & code, arg \\
invburr & reciprocal of burr & inverse Weibull; power &
  face comparison; with the scale held at $s$: all $x_i \le s$ & code, arg \\
betapr & shape2 $\to\infty$ / shape1 $\to\infty$ & gamma / inverse gamma &
  face comparison; none with the scale held & code, arg \\
\bottomrule
\end{tabular}
\end{center}

\paragraph{Free scale (\code{df32de2}, \code{11a5129}).} burr, invburr and
betapr used to hold the scale at $1$; they now estimate it, and a free scale
opens new faces. Burr XII, $S(x) = \{1 + (x/\sigma)^c\}^{-k}$: as
$k\to\infty$ with $\sigma k^{-1/c} = \lambda$ fixed,
$S(x) \to \exp\{-(x/\lambda)^c\}$, the Weibull; as $c\to\infty$ with $ck$
fixed, the Pareto with minimum $\sigma$ \cite{shao2004}. invburr is the
reciprocal, so its faces are the inverse Weibull and the power distribution.
Beta prime is $X = \sigma G_a/G_b$ with independent gamma variables: as
$b\to\infty$ with $\sigma/b$ fixed it tends to a gamma law, as $a\to\infty$
with $\sigma a$ fixed to an inverse gamma. Both shapes $\to\infty$ is
\emph{not} a lognormal limit: $\operatorname{Var}\log X = \psi'(a) + \psi'(b)
\to 0$. Each fitter compares its multistart interior fit with the package's
own fits of the limit models (faces win ties at $10^{-8}n$). In small
samples escapes are common (burr: 3 of 6 samples of size 20; betapr: 5 of
12). With the scale held (\code{mlburr(x, scale = 1)} is the 1.x fit) the
analysis below applies with $1$ replaced by the held scale; for betapr the
gamma and inverse-gamma faces need the scale to move, so the fixed-scale
MLE exists for every nonconstant sample.

\paragraph{burr with the scale held at 1.} $S(x) = (1+x^c)^{-k}$. With
$ck = \alpha$ fixed and $c\to\infty$, $S(x) = \exp\{-k\log(1+x^c)\} \to
x^{-\alpha}$ for $x>1$ and $S(x)\to 1$ for $x<1$: the Pareto law with
minimum $1$. All $x_i \ge 1$ is sufficient for escape, and it is also
necessary for this limit. Write $l_i = \log x_i$ and
$N = -\sum_i \min(l_i, 0)$. Along $k = \alpha/c$,
\[
  \log L = n\log\alpha - \textstyle\sum_i l_i - cN - \alpha P - Z\log 2 + o(1),
\]
with $P = \sum_i \max(l_i,0)$ and $Z$ the number of $x_i = 1$, so a single
$x_i < 1$ gives $N > 0$ and $\log L \to -\infty$. Profiling $k$ does not
change this. Mostly-Pareto samples therefore have a finite MLE, often with
large $c$ (e.g.\ $\hat c = 97.4$ for 199 points $\ge 1$ and one at $0.99$);
the fitter now finds it via a profile-likelihood start (commit
\code{9290aac}).

\subsection{Lomax and inverse Pareto: when the MLE exists}\label{sec:lomax}

% Adapted from an unpublished draft by J. Moss, "When does the Lomax maximum
% likelihood estimator exist?" (October 2026). Cleared for publication by the
% author on 2026-11-01.

The Lomax fitter stops with ``use \code{mlexp}'' when
$\overline{x^2} < 2\bar x^2$, that is, when the sample coefficient of
variation is below one. That test is only local: it makes the exponential a
strict local maximum of the profile, and a higher interior maximum can
exist. Del Castillo and Daoudi \cite{delcastillo2009} proved that the MLE exists when the
coefficient of variation exceeds one and conjectured (p.~687) that this is
also necessary; Renext's \code{flomax} \cite{renext} refuses to fit
otherwise. The conjecture is false. Grimshaw \cite{grimshaw1993} observed that the
generalized Pareto profile score can have two zeros on the heavy-tailed side,
and his algorithm compares each with the boundary. Below is an exact
criterion on a bounded interval.

\paragraph{Profile.} The density is $f(x) = \lambda\kappa(1+\lambda
x)^{-(\kappa+1)}$, $x > 0$. As $\lambda\to0$ with $\lambda\kappa\to\theta$ it
tends to the exponential with rate $\theta$, pointwise and hence in Hellinger
distance (Scheff\'e). Let $x\in(0,\infty)^n$, put $y_i = x_i/\bar x$,
$c^2 = n^{-1}\sum_i (y_i-1)^2$ and $\gamma = n^{-1}\sum_i (y_i-1)^3$ (divisor
$n$; $c<1$ is the same as $\overline{x^2}<2\bar x^2$). For fixed $\lambda$
the log-likelihood is strictly concave in $\kappa$ with maximiser
$\hat\kappa(\lambda) = n/\sum_i\log(1+\lambda x_i)$. With $q = \lambda\bar x$
and $A(q) = n^{-1}\sum_i\log(1+q y_i)$, $\hat\kappa = 1/A(q)$ and the profile
satisfies
\begin{equation}
  \ell_p(\lambda) - \ell_E = n\{\log q - \log A(q) - A(q)\},
  \qquad \ell_E = -n\log\bar x - n ,
  \label{eq:lomax-profile}
\end{equation}
where $\ell_E$ is the maximised exponential log-likelihood. The exponential
face is $q\to0$.

\begin{lemma}[\cite{delcastillo2009}, Proposition~2, proof of Corollary~1, Theorem~1]\label{lem:lomax-profile}
$\ell_p$ is continuous on $(0,\infty)$; (a) $\ell_p(\lambda)\to\ell_E$ as
$\lambda\to0$; (b) $\ell_p(\lambda)\to-\infty$ as $\lambda\to\infty$; (c)
$\ell_p(\lambda) - \ell_E = n(c^2-1)q/2 + O(q^2)$ as $q\to0$.
\end{lemma}

\begin{proof}
Since $n^{-1}\sum_i y_i = 1$ and $n^{-1}\sum_i y_i^2 = 1+c^2$,
$A(q) = q - (1+c^2)q^2/2 + O(q^3)$, so $A(q)/q$ is analytic near $0$ with
value $1$ there and $\ell_p - \ell_E = -n\{\log(A(q)/q) + A(q)\}$ has the
stated expansion. This gives (a) and (c). For (b), $q\to\infty$ and
$A(q) = \log q + n^{-1}\sum_i\log y_i + o(1)$, so $\log q - A(q)$ stays
bounded while $\log A(q)\to\infty$.
\end{proof}

By (a) and (b) the MLE exists if and only if $\ell_p(\lambda)\ge\ell_E$ for
some $\lambda$. By (c), $c>1$ suffices, and $c<1$ makes the exponential a
strict \emph{local} maximum. Define, for $u\ge0$,
\[
  \psi(u) = \frac1n\sum_{i=1}^n\log\bigl(e^{-u} + u y_i\bigr),
  \qquad
  \bar u = \frac{\bar x}{(\prod_i x_i)^{1/n}} \ge 1 ,
\]
and let $W$ be the principal branch of the Lambert $W$ function
\cite{corless1996}, the inverse of $u\mapsto ue^u$ on $[0,\infty)$.

\begin{theorem}\label{thm:lomax}
Let $x\in(0,\infty)^n$.
(a) For every $\lambda>0$, $\ell_p(\lambda)-\ell_E$ has the sign of
$-\psi(W(\lambda\bar x))$.
(b) The MLE exists if and only if $\psi(u)\le0$ for some $u>0$.
(c) $\psi(u)>0$ for all $u\ge\bar u$, so every maximiser satisfies
$\hat\lambda\bar x < \bar u e^{\bar u}$.
\end{theorem}

\begin{proof}
(a) Put $q = \lambda\bar x$ and $u = W(q)$, so $q = ue^u$. By
\eqref{eq:lomax-profile}, $\ell_p\ge\ell_E$ if and only if
$q\ge A(q)e^{A(q)}$, and since $z\mapsto ze^z$ is strictly increasing on
$[0,\infty)$, if and only if $u\ge A(q)$; the same holds with strict
inequalities. Finally
$u - A(q) = u - n^{-1}\sum_i\log(1+ue^uy_i) = -\psi(u)$.
(b) By the lemma and (a), since $\lambda\mapsto W(\lambda\bar x)$ maps
$(0,\infty)$ onto $(0,\infty)$.
(c) $\psi(u) > n^{-1}\sum_i\log(uy_i) = \log u - \log\bar u\ge0$ for
$u\ge\bar u$. A maximiser has $\ell_p(\hat\lambda)\ge\ell_E$, so
$\psi(W(\hat\lambda\bar x))\le0$ and $W(\hat\lambda\bar x)<\bar u$.
\end{proof}

The function $\psi$ depends on the sample only through the empirical distribution of the
$y_i$, so the answer is invariant under rescaling and replication. Since
$y\mapsto\log(e^{-u}+uy)$ is concave, a mean-preserving spread of the $y_i$
lowers $\psi$.

\begin{proposition}\label{prop:lomax-expansion}
As $u\to0$, $\psi(u) = \tfrac12(1-c^2)u^2 + \tfrac16(2\gamma-1)u^3 + O(u^4)$.
Hence the MLE exists if $c>1$, or if $c=1$ and $\gamma<1/2$; in the second
case also on an open neighbourhood of $x$, which contains samples with $c<1$.
\end{proposition}

\begin{proof}
With $\delta_i = y_i - 1$, $e^{-u}+uy_i = 1+z_i$,
$z_i = u\delta_i + u^2/2 - u^3/6 + O(u^4)$, and
$\log(1+z_i) = u\delta_i + \tfrac12(1-\delta_i^2)u^2
+ (\tfrac13\delta_i^3 - \tfrac12\delta_i - \tfrac16)u^3 + O(u^4)$. Average
over $i$, using $\sum_i\delta_i = 0$. The neighbourhood statement follows
from continuity of $\psi(u_0)$ in $x$ at a $u_0$ with $\psi(u_0)<0$, and
because $c$ is continuous and not locally constant.
\end{proof}

\begin{proposition}\label{prop:lomax-counter}
Let $n\ge2$, $k = \lfloor n/2\rfloor$, and let $x_\varepsilon$ have $k$
copies of $\varepsilon\in(0,1)$ and $n-k$ copies of $1$. Then $c<1$ for all
$\varepsilon$, and the MLE exists for all sufficiently small $\varepsilon$.
The samples with $c<1$ and an MLE contain a nonempty open set, which has
positive probability under every exponential and Lomax law.
\end{proposition}

\begin{proof}
With $w = k/n\le\frac12$,
$c^2 = w(1-w)(1-\varepsilon)^2/(1-w+w\varepsilon)^2 < 1$, since
$w(1-w)\le(1-w)^2$ and $(1-w)(1-\varepsilon) < 1-w+w\varepsilon$. As
$\varepsilon\to0$, $\psi(u)\to -wu + (1-w)\log\{e^{-u} + u/(1-w)\}$, which
tends to $-\infty$ as $u\to\infty$; take $u_0$ where the limit is negative
and apply Theorem~\ref{thm:lomax}(b). Openness follows from continuity of $c$
and $\psi(u_0)$ in $x$.
\end{proof}

\paragraph{Examples} (checked numerically against
\code{extraDistr::dlomax}; gap $=\ell_p(\hat\lambda)-\ell_E$).

\begin{center}
\small
\begin{tabular}{@{}lrrrrr@{}}
\toprule
sample & $c$ & $\min\psi$ & gap & $\hat\lambda$ & $\hat\kappa$ \\
\midrule
$(1, 100)$ & 0.9802 & $-0.214$ & 0.5901 & 0.7722 & 0.4055 \\
$(0.001, 4, 6)$ & 0.7481 & $-0.386$ & 1.3383 & 1394.1 & 0.1619 \\
$(0.001, 1, 2)$ & 0.8158 & $-0.183$ & 0.6524 & 1039.3 & 0.1961 \\
$(1, 29)$ & 0.9333 & $-5.1\times10^{-4}$ & 0.0019 & 0.2370 & 0.8787 \\
$(1, 1, 1, 12, 12)$ & 0.9979 & $-5.5\times10^{-6}$ & $1.9\times10^{-4}$ & 0.0360 & 6.0716 \\
$(1, 10)$ & 0.8182 & $0$ ($u\to0$) & no MLE & & \\
\bottomrule
\end{tabular}
\end{center}

$(1,1,1,12,12)$ lies near the case $c = 1$, $\gamma < 1/2$ of
Proposition~\ref{prop:lomax-expansion} ($\gamma = 0.41$). For a pair, $\psi$ increases
with the ratio $r$ of the smaller to the larger observation (a smaller ratio
is a mean-preserving spread), so a pair has an MLE exactly when
$r<\varepsilon^\ast = 0.034711$ (larger/smaller $> 28.8$); for two
independent exponential observations this has probability
$2\varepsilon^\ast/(1+\varepsilon^\ast) = 6.7\%$. In simulations with
$10^4$ samples per setting the moment rule was wrong for 7.0\% of
exponential pairs and up to 31\% of Lomax pairs ($\kappa = 0.5$), for at most
1.1\% of samples of size 3 and 5, and for none of size 20 or 50, although
63\% of exponential samples of size 50 have $c<1$.

\paragraph{Consequences for the fitter} (implemented 2026-10-29 in
\code{lomax/mle.hpp}: grid in $\log u$ plus refinement of every local
maximum, compressed scan for $n>4096$). (i) Existence is decided by the sign
of $\min\psi$ on the bounded interval $(0,\bar u]$; the minimisation is a
smooth one-dimensional global search, done numerically. (ii) If the minimum
is nonpositive, every maximiser has $u = W(\hat\lambda\bar x)$ in
$\{u\in(0,\bar u):\psi(u)\le0\}$, so the profile only needs to be maximised
over $\lambda\in(0,\bar u e^{\bar u}/\bar x)$; it can have more than one
interior maximum, so this search is global too. (iii) The package's tie
convention (faces win ties at $10^{-8}n$) sends attained ties ($\min\psi = 0$)
and gaps below $10^{-8}n$ to the exponential escape. (iv) $c>1$ needs no search for existence.

\paragraph{Inverse Pareto.} If $X$ has the inverse Pareto law with shape
$\tau$ and scale $\theta$ (\code{actuar::dinvpareto}), then $1/X$ is Lomax
with $\kappa=\tau$, $\lambda=\theta$. The log-likelihoods of $x$ and $1/x$
differ by $-2\sum_i\log x_i$, free of parameters, and the same holds for the
inverse exponential against the exponential. So everything above applies to
$1/x_i$, with limit model \code{mlinvexp}.

\subsection{Escape, formerly silent; now detected by comparing faces}

\begin{center}
\small
\begin{tabular}{@{}l p{3.4cm} p{3.2cm} p{5.4cm} l@{}}
\toprule
id & escape & limit model & probe result (before the fix) & firm \\
\midrule
std & $\nu\to\infty$ & normal &
  normal data, $n=200$: $\hat\nu = 54874$, $\ell$ $2\times10^{-4}$ below
  \code{mlnorm}; at $n=30$ $\hat\nu = 9$ (finite) & arg, probe \\
std & $\nu\to 2^+$ & location-scale $t_2$ (infinite variance) &
  not representable in the mean/sd parameterisation & arg \\
snorm & $\xi\to 0$ or $\infty$ & translated half-normal &
  exponential data: $\hat\xi = 7.4\times10^{10}$ & arg, probe \\
sstd & $\nu\to\infty$; $\xi\to0,\infty$; $\nu\to2^+$ & snorm; half-$t$; two-piece $t_2$ &
  normal data: $\hat\nu = 343$, $\ell$ $0.075$ \emph{below} \code{mlsnorm};
  on $n=20$ family data the oracle finds $\xi\to\infty$ with higher $\ell$ & probe \\
\bottomrule
\end{tabular}
\end{center}

The skew families use the Fern\'andez--Steel construction \cite{fs1998}; as
$\xi\to\infty$ (or $0$) all mass moves to one side of the mode, which gives
half-distributions whose endpoint (the mode) is estimated, not fixed at $0$.

\paragraph{std.} Let $K = n^{-1}\sum_i z_i^4$ be the sample kurtosis, with
$z_i$ standardised by the normal MLE. Expanding the unit-variance $t$
density in $\eta = 1/\nu$ gives the inward score of the profile
log-likelihood at the normal boundary,
\[
  \ell_{\mathrm{profile}}'(0^+) = \frac{n(K-3)}{4}.
\]
So $K < 3$ makes the normal a strict \emph{local} maximum and $K \le 3$ is
necessary for the normal to be the global supremum. It is not sufficient: a
two-cluster sample with $K = 2.33$ has its supremum at the $\nu\to2^+$ face
instead. The sd parameterisation (fGarch's \code{dstd}) cannot represent
that face at all. Detection must compare the boundary faces with the
interior, not threshold $\hat\nu$.

Fixed (\code{3e9ffee}): the fitter compares the interior BFGS fit (plus a
restart near large $\nu$ when $K < 3$) with the closed-form normal fit and
the best location-scale $t_2$ fit, and reports \code{NormalBoundary}
(``use \code{mlnorm}'') or \code{DfTwoBoundary}. Two old test samples turned
out to have their supremum on the $t_2$ face.

\paragraph{snorm: an exact rule.}\label{sec:snorm} Write the two-piece normal
with mode $c$ and left/right scales $a, b$. For fixed $c$ let
$A(c) = \sum_{x_i<c}(x_i-c)^2$, $B(c) = \sum_{x_i \ge c}(x_i-c)^2$. The
scales profile out in closed form ($a/b = (A/B)^{1/3}$), giving
\[
  \ell_{\mathrm{profile}}(c) = n\log\sqrt{2/\pi} - \frac n2 \log\frac{Q(c)}{n}
  - \frac n2, \qquad Q(c) = \bigl(A(c)^{1/3} + B(c)^{1/3}\bigr)^3 .
\]
Hence the snorm MLE is a one-dimensional minimisation of $Q$ over
$[x_{(1)}, x_{(n)}]$, and the MLE escapes exactly when $Q$ is minimised only
at an endpoint ($x_{(1)}$: right half-normal; $x_{(n)}$: left half-normal).
On each interval between order statistics the stationarity condition is a
polynomial of degree at most seven, so all candidates can be enumerated.
Both half-normal faces are always strict local maxima, which explains why a
local optimiser gets stuck there. (Derivation checked; see the report.)

Implemented (\code{93c7422}): BFGS is replaced by this exact algorithm
(all stationary points on every interval via root isolation,
$O(n\log n)$); the faces are reported as \code{RightHalfNormalBoundary} /
\code{LeftHalfNormalBoundary}.

\paragraph{sstd.} No cheap global rule is known. The inward score at the
snorm face is again $(\sum_i w_i^4 - 3n)/4$; it is a local filter only. In
mode/scale coordinates $(c, a, b, \nu)$ with $\xi = \sqrt{b/a}$ and the
ordinary $t_\nu$, the density is $2\,t_\nu((x-c)/s)/(a+b)$, $s = a$ left of
$c$ and $b$ right of it. The closure has three faces: $\nu\to\infty$ (snorm,
with its half-normal corners), $\xi\to\infty$ or $0$ (a translated half-$t$
with endpoint at $x_{(1)}$ or $x_{(n)}$ and $\nu \in [2,\infty]$), and
$\nu\to2$ (two-piece $t_2$, infinite variance).

Implemented (\code{73f4fd5}): the fitter compares the best of seven interior
BFGS starts with each face, fitted directly: snorm by the exact $Q$-profile;
the half-$t$ by profiling the scale (its score
$\sum_i(\nu+1)w_i^2/(\nu+w_i^2) - n$ decreases in the scale, so the root is
unique) and a grid-plus-golden-section search in $\eta = 1/\nu \in
[0,\tfrac12]$; the $t_2$ face by BFGS from quantile starts and from every
interior candidate. Faces win ties within $10^{-8}n$: an interior iterate
running off towards a face approaches the face value from below. Statuses
\code{SnormBoundary} (``use \code{mlsnorm}''),
\code{Right/LeftHalfNormalBoundary} (``use \code{mlhnorm}'', on $-x$ for the
left corner), \code{Right/LeftHalfTBoundary} (names the df) and
\code{DfTwoBoundary}.

Checked against a slow oracle (multistart \code{optim} on
\code{fGarch::dsstd}, plus each face fitted directly) on 45 samples: every
interior fit matches the oracle, and every reported face is the best face
with the oracle's interior supremum not above it. Faces are common in small
samples: of 30 samples of size 20 from skew-$t$ laws with $\nu = 5$ or $10$,
23 have no interior maximum (14 snorm or half-normal, 8 half-$t$, 1 $t_2$).
Normal data go to the snorm face, exponential data to the half-$t$ face
($\hat\nu\approx3.3$, heavier than the half-normal tail), skewed Cauchy data
to the $t_2$ face. The old ``stops at $\hat\nu = 343$'' was an overflow, not
an optimum: see \S\ref{sec:bugs}.

\subsection{No MLE at all: ged and sged}

For the exponential-power (GED) family with free location, scale and shape
$p$, put the location at an observation and profile the scale. Then
\[
  \frac{\ell_p}{n} = -\frac{\log q}{p} + \tfrac12\log p + O(1)
  \to \infty \quad (p\to 0),
\]
where $q = (n-k)/n$ and $k$ is the number of observations at that location.
The likelihood is unbounded for \emph{every} nonconstant sample
\cite{tumlinson2015}; sged contains ged ($\xi = 1$), so the same holds.
Numerical check (normal data, $n = 40$, location at $x_1$): the profile
log-likelihood is $-48.4$ at $p = 2$, $-65.3$ at $p = 0.03$ and $-19.0$ at
$p = 0.01$, above the normal MLE's $-46.0$, and rising.

The $\nu\to\infty$ limit (uniform on $[\mu\pm\sqrt3\,\sigma]$; for sged an
ordinary uniform, the skew becomes unidentified) is a local maximum only,
with inward score $-\infty$ at the boundary. A well-posed fit needs a lower
bound on the shape, which is a model change, not a numerical tolerance.

Decided and implemented (D6, \code{97bc3a6}): both fit on $\nu \ge 1$, the
log-concave GED, which contains the Laplace and every lighter tail. The
supremum is then at $\nu = 1$ (attained, the Laplace), at a finite
$\nu > 1$, or at $\nu \to \infty$ (the uniform; status
\code{UniformBoundary}, ``use \code{mlunif}''); sged also compares the
half-GED faces at $x_{(1)}$ and $x_{(n)}$. The behaviour sweep's unrestricted
oracle still beats the restricted fit when its optimum has $\nu < 1$ (one
outlier sample); that is the restriction, not a fitter failure.

\subsection{Escape, fixed 2026-10-28}

\paragraph{invburr with the scale held at 1.} $F(x) =
\bigl(x^\gamma/(1+x^\gamma)\bigr)^\tau$, and $X \sim
\text{invburr}(\tau,\gamma)$ iff $1/X \sim \text{burr}(c=\gamma,k=\tau)$.
With $\gamma\tau = \alpha$ fixed and $\gamma\to\infty$, $F(x)\to x^\alpha$
on $(0,1]$ and $F(x)\to 1$ for $x>1$: the power distribution on $(0,1]$,
the image of the burr limit under $x\mapsto 1/x$. All $x_i \le 1$ is
necessary and sufficient (reciprocal of the burr argument). It used to give
``Inverse Burr MLE failed''; it now returns status \code{AllLeOne} and the
message names \code{mlpower} (commit \code{411ea90}).

\subsection{Limit models added to the package (D7)}

The escape messages need a model to point to. Two were added
(\code{dev/research/ml-behaviour-limit-models.md}):
\begin{itemize}
  \item hnorm, the translated half-normal (\code{8807e8b}): the boundary MLE
    $\hat\mu = x_{(1)}$, $\hat\sigma^2 = \overline{(x-\hat\mu)^2}$. The R
    package defines no densities of its own, so \code{dml}, \code{pml},
    \code{qml}, \code{rml} call \code{extraDistr::dtnorm} with
    $(\text{mean}, \text{sd}, a) = (\mu, \sigma, \mu)$ through an argument
    map in the README front matter (\code{cb8b9ac}). \code{dtnorm} excludes
    its lower endpoint, so the R fitter reports $\hat\mu$ as the largest
    double below $x_{(1)}$; then the log-likelihood equals the sum of
    \code{dml}.
  \item lst, the location-scale $t$ with $\mathrm{df} \ge 1$
    (\code{92eef3f}): the target of std's $\nu\to2^+$ face. Its own
    $\mathrm{df}\to\infty$ face stops with ``use \code{mlnorm}''.
\end{itemize}
The half-$t$ and two-piece $t_2$ faces of sstd and the half-GED faces of
sged still have no package model; their messages name the limit.

\subsection{Not closed, but the MLE always exists}\label{sec:exists}

\paragraph{invgauss.} As the mean $\mu\to\infty$ with shape fixed, the
inverse Gaussian converges to the L\'evy distribution. The MLE
$\hat\mu = \bar x$ is always finite, so there is no escape. Probe on L\'evy
data: $\hat\mu = 1663$, finite but meaningless.

\subsection{MLE on the boundary of the parameter space}

These are not escapes (the boundary point is in the family) but deserve a
line of documentation.
\begin{itemize}
  \item zip: $\hat\pi = 0$ (the Poisson fit) when the data have fewer zeros
    than the Poisson predicts. Probe agrees with \code{mlpois}.
  \item zipf: $\hat s < 0$ is rejected with status \code{ShapeNegative};
    the boundary $s = 0$ is the uniform law on $\{1,\dots,N\}$.
  \item naka: the probe returned $\hat m = 0.275 < 1/2$. The density is
    integrable for every $m > 0$ ($X^2$ is gamma), and the R packages
    \code{nakagami} and \code{VGAM} and scipy allow $m > 0$; $m \ge 1/2$ is
    a radio-fading convention. If $m \ge 1/2$ were imposed, $m = 1/2$ (the
    half-normal) is an attained boundary, not an escape (decision D5).
  \item unif, dunif, pareto, power: the MLE sits at an order statistic.
\end{itemize}

\subsection{Closed apart from point masses}

norm, lnorm, logitnorm, exp, rayleigh, pois, geom, lgser, gamma, lgamma,
invgamma, weibull, invweibull, gumbel, logis, llogis, laplace, cauchy, beta,
paralogis (arg, not probed), and kumar. For kumar, with $t_i = -\log x_i$ and
$S(a) = \sum_i \log(1 - e^{-a t_i})$, the shape $b = -n/S(a)$ profiles out
and the profile log-likelihood tends to $-\infty$ both as $a\to0$ (like
$-n\log\lvert\log a\rvert$) and as $a\to\infty$ (like
$a(n t_{\min} - \sum_i t_i)$), so the MLE exists for every nonconstant
sample; the old failure on data piled near $0$ and $1$ was numerical
(fixed, \code{808ddf5}).

\paragraph{Fixed degrees of freedom.} With $\nu$ held (\code{mlstd(x, nu =
v)}, \code{mllst(x, df = v)}) the normal and $t_2$ faces disappear. A value
repeated $k$ times still matters: putting the location there and shrinking
the scale gives $L \propto \text{scale}^{\,\nu(n-k)-k}$, so the likelihood is
unbounded when $\nu(n-k) < k$; the fitter reports this, and stops with a
diagnostic in the equality case.

\paragraph{fatigue: why the location is fixed.} The fitter estimates
$(\alpha,\beta)$ with $\mu = 0$ fixed. With $\mu$ free the likelihood is
unbounded for every sample with a unique minimum, so the three-parameter
MLE never exists. Write $t = x - \mu$ and
\[
  \log f(t) = \log(t+\beta) - \tfrac32\log t - \tfrac12\log\beta - \log(2\alpha)
  - \tfrac12\log(2\pi) - \frac{t/\beta + \beta/t - 2}{2\alpha^2}.
\]
Take $\beta = \varepsilon$, $\alpha^2 = s/\varepsilon$ and
$\mu = x_{(1)} - c\varepsilon^2/s$ with $s, c > 0$ fixed. For the smallest
observation, $t_1 = c\varepsilon^2/s$ and the $\log\varepsilon$ terms add to
$(1 - 3 - \tfrac12 + \tfrac12)\log\varepsilon$, while the exponent tends to
$-1/(2c)$; hence $\log f(t_1) = -2\log\varepsilon + O(1)$. For the others,
$t_i \to d_i = x_i - x_{(1)} > 0$ and
$\log f(t_i) \to \log\tfrac12 - \tfrac12\log(2\pi s d_i) - d_i/(2s)$, the
log-density of half of a scaled $\chi^2_1$ law: as $\alpha\to\infty$ with
$\alpha^2\beta$ fixed, half of the Birnbaum--Saunders mass collapses onto
the endpoint. So $\log L = -2\log\varepsilon + O(1) \to \infty$ (numerical
check: $+4.6 = 2\ln 10$ per decade of $\varepsilon$; a test in
\code{fatigue/test.R}). This is the same defect as the three-parameter
lognormal; the two-parameter fit with a known threshold is the usable model.

\subsection{Families added 2026-10-29}\label{sec:new-families}

Thirty-one further families, implemented between the 2026-10-29
parameter-rules design and its P0--P3 jobs. Existence claims below are as
classified in the ten P0 group notes and the per-family implementation
reports; ``proved'' means a direct argument in one of those notes, ``probe''
means a numerical multistart/grid search (evidence, not a certificate of a
global maximum), and ``conj'' means an open classification. Every P0 job
flags the same residual limitation: multistart optimisation cannot certify a
global maximum, so a probed existence boundary is always provisional.

\subsubsection*{Discrete and count families}

\begin{center}
\small
\begin{tabular}{@{}l p{4.2cm} p{8.2cm}@{}}
\toprule
id & closure & existence, boundary and multimodality \\
\midrule
bnbinom & NB ($\alpha,\beta\to\infty$ fixed ratio, or the mirror swap
  size$\leftrightarrow\beta$); Poisson ($r,\beta,\alpha\to\infty$,
  $r\beta/\alpha$ fixed); point mass at 0 (any of $r\to0$, $\beta\to0$,
  $\alpha\to\infty$, all-zero data) &
  Proved for samples with a positive count when $\alpha$, or two
  coefficients, are held (compact parameter space). For the fully free fit,
  exists iff some finite parameter point beats every reachable NB/Poisson
  face, an implicit criterion, not a closed-form statistic; whether every
  sample with $\overline{x^2}-\bar x^2\le\bar x$ selects Poisson is open
  \cite{irwin1975,xekalaki1984} (probe/conj) \\
cmp & adjacent two-point law ($\nu\to\infty$); point mass (constant data,
  all free) &
  Own d/p/q/r. Interior MLE exists and is unique iff $b(\bar x)<T<g(\bar x)$,
  $T=n^{-1}\sum_i\log(x_i!)$, $g,b$ the $\nu=0$ and $\nu\to\infty$ limits of
  $\E\log(X!)$ at fixed mean (proved: the $\eta$-score is strictly monotone
  in $\nu$ by a variance/covariance argument). $\lambda=0$ and ($\nu=0$,
  $\lambda<1$) are attained R7 members since 2026-10-29, not escapes
  (proved, code) \\
cmb \cite{kadane2016} & interior/adjacent/endpoint two-point laws as
  $\nu\to\pm\infty$ &
  $\mathrm{prob}\in\{0,1\}$ and $\mathrm{size}=0$ attained (R7, code). Free
  size unimplemented: whether the reference's $10^7$ cap is a model bound or
  an allocation limit is undecided (conj) \\
zeta, yulesimon & point mass at 1 as shape $\to\infty$ &
  Existence iff some $x_i>1$ (proved: zeta's shape score is model mean minus
  empirical mean of $\log X$, strictly decreasing; Yule--Simon's is
  $n-\sum_i x_i$) (proved) \\
tpois & point mass at the free support endpoint opposite a held bound, when
  data sit at that extreme &
  Closed-form conditional-mean solve on any data-compatible bracket
  $[a,b]$; R7a (\S\ref{sec:holds-restrictions}): the lower bound is
  $\max(0,\lfloor a\rfloor+1)$, never an atom at $a$ itself. Singleton
  support makes every $\lambda>0$ an unidentified maximizer; still rejected
  by the fitter (B7, open) \\
tbinom & point mass at a bound as $\mathrm{prob}\to0,1$ or
  $\mathrm{size}\to\infty$ at fixed interior prob &
  Closed for held size (all six proper subsets, proved). Free size with
  free probability needs a compactification proof not yet supplied; four
  hold subsets stay \code{not-implemented} for that reason (conj) \\
dlaplace & point mass at location as scale $\to0$ &
  Median location (possibly nonunique), closed-form scale
  $d/(\sqrt{1+d^2}+1)$, $d=$ mean absolute deviation (proved). R7: restricted
  to integer location although the reference accepts real location, because
  the reference's real-argument values do not sum to one on the integer
  lattice \\
nhyper & none: the known-$N$ design makes the likelihood bounded ($\le1$)
  on a finite integer grid &
  Log-concave in $m$ (adjacent-ratio argument), so bisection finds the
  (possibly tied) global maximum; no escape is possible under this design
  (proved). As $m,n\to\infty$ with $m/(m+n)\to p$, $T-r\to\mathrm{NB}(r,p)$;
  the correct referral is \code{mlnbinom(x - r, size = r)}, not
  \code{mlnbinom(x, size = r)} \cite{berrymielke1998} \\
dnorm, dgamma & point mass (range 0) or shifted \code{mlbinom(size=1)}
  (range 1) &
  Existence iff range $\ge2$: compact likelihood level sets from a uniform
  per-cell probability bound plus grouped-normal concavity, citing
  Burridge \cite{burridge1981} for dnorm (proved); dgamma: analogous
  compactness with a different per-cell bound; point mass at zero as
  $\mathrm{rate}\to\infty$ (shape held) or $\mathrm{shape}\to0$ (rate held) \\
skellam & none: zero-rate faces are attained R7 members, not escapes &
  Joint MLE matches the empirical mean and variance exactly (closed score
  equations, proved). Compactness (each mass is bounded by a Stirling bound
  on the larger-rate Poisson factor, $\to0$ as either rate diverges) gives
  an attained finite MLE for \emph{every} integer sample (proved) \\
\bottomrule
\end{tabular}
\end{center}

\subsubsection*{Half-threshold and real-line shape families}

\begin{center}
\small
\begin{tabular}{@{}l p{4.2cm} p{8.2cm}@{}}
\toprule
id & closure & existence, boundary and multimodality \\
\midrule
hnorm (held $\mu$) & translated half-normal, $\mu$ kind \code{threshold} &
  With $\mu$ free, always exists if nonconstant: the profile increases
  $\mu$ to $\min(x)$, with closed-form $\hat\sigma^2$ (proved). R7a
  (2026-10-28) makes the endpoint exactly attained, replacing the earlier
  ``largest double below $x_{(1)}$'' convention \\
hcauchy & translated endpoint $\mu$ (required threshold) &
  Scale score $n-\sum_i 2\sigma^2/(\sigma^2+x_i^2)$ strictly decreasing from
  $n-2k$ to $-n$, $k=\#$ zeros: unique MLE iff $2k<n$; unbounded if $2k>n$;
  finite unattained supremum at $2k=n$ (proved) \\
ht & half-normal escape, $\nu\to\infty$; restricted $\nu\ge1$ (collapse,
  $c=\nu$, Prop.~\ref{prop:collapse}) &
  Scale score decreasing from $\nu(n-k)-k$ to $-n$: unique root iff
  $\nu(n-k)>k$ (proved). Multiple local maxima occur even in small samples
  (a five-point \code{README} example has maxima at $\hat\nu=0.415$ and
  $1.17$, both above the half-normal face) (probe) \\
foldnorm & none (compact level sets for every nonconstant sample) &
  $f\le2\times$ the parent normal density gives attained MLE for every
  nonconstant sample; constant data collapse (proved). Multimodality proved
  to occur: a systematic 20{,}000-sample scan found one sample with two
  interior maxima (3 stationary points); a moment-only classifier
  ($m_4$ vs.\ $3m_2^2$) is only local, with an explicit counterexample
  (99 ones, one five) \\
huber (free $\epsilon$) & normal ($\epsilon\to\infty$); Laplace
  ($\epsilon\to0$, $\sigma$ free, face scale $b=\sigma/\epsilon$, not
  $\sigma\epsilon$) &
  For $n=2$ distinct points the fully free MLE never exists: both faces
  dominate every finite $\epsilon$ (proved exactly). M4 simulation: normal
  data escape to \code{mlnorm} $60$--$72\%$ of the time
  ($n=20,100,1000$), Cauchy data up to $62\%$ at $n=1000$; kept kind
  \code{design}, required (A5, probe) \\
sn & translated half-normal, either side, $\alpha\to\pm\infty$ (with a
  $\log2$ penalty per observation tied exactly at a held $\xi$) &
  Normal ($\alpha=0$) is \emph{always} a critical point of the profile for
  nonconstant data (singular Fisher information there, proved), so a local
  search alone cannot certify the global fit. Escape-frequency Monte Carlo
  at $\alpha=5,10$, $n=50$: $0.34,0.65$ here against published $0.139,0.345$
  \cite{azzalini1985,sartori2006}; unresolved (open) \\
emg & normal ($\lambda\to\infty$); translated exponential ($\sigma\to0$,
  endpoint $=\min(x)$, rate $=n/\sum(x-\min)$) &
  Interior stationary point satisfies $\hat\mu+1/\hat\lambda=\bar x$
  (diagnostic only). Local trap proved concretely on $x=-2{:}2$: normal is
  a strict local maximum ($m_4-3m_2^2<0$) while the exponential face has
  strictly higher likelihood (proved) \\
slash & point-mass collapse at a tied observation, $\sigma^{n-2k}$
  ($c=1$ automatically, Cauchy-order tails) &
  Attained for $k<n/2$ (proved). At $k=n/2$ exactly, a variational
  criterion compares the interior profile to the exact half-mass value
  $B_v$, not reducible to a closed sample statistic. Multimodality proved
  for an explicit 20-point sample, certified with exact rational-arithmetic
  inequalities \\
bhatt & normal ($a\to0$, now R7-attained, not an escape); uniform
  ($\sigma\to0$, attained only if range $>2a$ with $a$ held) &
  $f\le1/(2a)$ and $f\le1/(\sqrt{2\pi}\sigma)$ exclude joint divergence
  except at constant data (proved). Escape-frequency simulation (100
  samples/cell): at $n=20$, normal-generated data pick the uniform face
  $73\%$ of the time and the normal face only $21\%$ (probe) \\
\bottomrule
\end{tabular}
\end{center}

\subsubsection*{Generalized beta and Pareto-type families}

\begin{center}
\small
\begin{tabular}{@{}l p{4.2cm} p{8.2cm}@{}}
\toprule
id & closure & existence, boundary and multimodality \\
\midrule
gb2 & six faces: generalized gamma, inverse GG, lognormal, Pareto, power,
  ``double Pareto'' (no package fitter) &
  Proved for 5 of 16 proper hold subsets (concave profiles). Escape-probe
  (7 generators, $n=50,500$, 20 reps): at $n=50$ even GB2-generated data
  escape 12--14/20 times; lognormal data split near 50/50 between the
  generalized-gamma and inverse-GG corner faces at $n=500$ (Self--Liang
  boundary-mixture behaviour) (probe) \\
genbeta & unit gamma ($\mathrm{shape1}\to\infty$, $\mathrm{shape3}\to0$,
  product finite); power ($\mathrm{shape1}\to0$, $\mathrm{shape3}\to\infty$),
  both at held scale &
  \code{scale} (upper endpoint) is \code{threshold}, required: free scale
  is \code{no-mle} for every shape subset (A4). $\mathrm{shape2}<1$ gives
  an upper-endpoint singularity of exponent $\mathrm{shape2}-1$,
  independent of the collapse mechanism, so a bare
  $\mathrm{shape2}\ge1$ restriction alone would not restore existence
  (\S\ref{sec:holds-restrictions}) \\
pareto2 & shifted Lomax; exponential face under the shifted
  Theorem~\ref{thm:lomax} criterion &
  With \code{min} held, data-dependent exactly as Lomax (proved). Free
  threshold is \code{no-mle} (A4). Threshold collapse exponent $c=$
  shape, with \emph{no} independent endpoint singularity (density finite
  at \code{min}) \\
pareto3 & shifted log-logistic at held \code{min} &
  Free threshold \code{no-mle}; collapse exponent $c=$ shape, \emph{plus}
  an independent fixed-scale endpoint singularity whenever shape $<1$: the
  counterexample to ``collapse $c\ge1$ suffices'' (\S\ref{sec:holds-restrictions}) \\
pareto4 & shifted Burr at held \code{min} (Weibull/Pareto faces) &
  Free threshold \code{no-mle}; collapse exponent
  $c=\mathrm{shape1}\cdot\mathrm{shape2}$, independent endpoint singularity
  when $\mathrm{shape2}<1$ \\
fpareto & shifted GB2 at held \code{min} (six faces) &
  Free threshold \code{no-mle}; collapse exponent
  $c=\mathrm{shape1}\cdot\mathrm{shape2}$, independent singularity when
  $\mathrm{shape2}\cdot\mathrm{shape3}<1$ \\
nsbeta & none (bounded support) &
  Both endpoints \code{threshold}: independent singularities at each end
  ($\mathrm{shape1}<1$ at \code{min}, $\mathrm{shape2}<1$ at \code{max});
  free-endpoint subsets all \code{no-mle} (A4) \\
\bottomrule
\end{tabular}
\end{center}

\subsubsection*{Extreme-value threshold and branch families}

\begin{center}
\small
\begin{tabular}{@{}l p{4.2cm} p{8.2cm}@{}}
\toprule
id & closure & existence, boundary and multimodality \\
\midrule
sgomp & none: $\eta=0$ (exponential) is an ordinary attained member,
  domain $\eta\ge0$ &
  Proved attained for every nonconstant sample (fixed-$b$ shape score
  strictly monotone with a bounded root; two-sided rate bounds for
  nonconstant data). Collapse only at constant/all-zero data with the
  relevant scale free \\
frechet & inverse Weibull on $x-\mu$ ($\mu$ held, any $\lambda$); Gumbel
  ($\lambda\to\infty$, free location) &
  $\lambda\ge1$ (collapse, $c=\lambda$, exempt when $\mu$ held) is exactly
  Prop.~\ref{prop:collapse} applied to Fr\'echet. A second, distinct spike
  ($\lambda\to0$ with $\mu\to\min$, present even at held $\sigma$) is not
  excluded by the $\lambda\ge1$ bound and needs $\lambda$ itself free \\
tlambda & logistic ($\lambda\to0$); uniform at $\lambda=1$ (attained
  alias) and $\lambda=2$ (alternate alias, held only) &
  Unrestricted likelihood is \emph{always} unbounded (location at an
  observation, $\sigma=\lambda A$, $A\to\infty$: $\ell=k\lambda\log2-n\log
  \lambda+O(1)\to\infty$ for every $\lambda>0$, proved). Two simultaneous
  restrictions on one shape, $\lambda\in[-1,1]$: collapse ($c=-1/\lambda$
  for $\lambda<0$, exempt when $\mu$ held) and branch ($\lambda\le1$, no
  exemption). Attained for $n\ge3$, every tie $k<n/2$ (explicit coercivity
  bound, proved); location likelihood proved \emph{not} unimodal (nearby
  local maxima as $\sigma\to0$ at $\lambda=-1$) \\
\bottomrule
\end{tabular}
\end{center}

\section{Existence under holds and restrictions}\label{sec:holds-restrictions}

The parameter-rules design (\code{dev/roadmap/2026-10-29-parameter-rules.md},
R1--R8, decided 2026-10-29) gives every coefficient a \emph{kind}:
\code{design} (fixed by data collection, or a tuning constant such as
huber's $\epsilon$; required), \code{threshold} (a shift of a family whose
standard form starts at a fixed point; required), \code{endpoint} (a support
bound; estimated by default) or \code{regular} (estimated by default). Rule
R5 restricts a \emph{shape} coefficient only when the family has no MLE for
almost any sample, through one of four mechanisms.

\begin{itemize}
  \item \emph{Collapse}: a location coincides with $k$ of $n$ observations
    and the likelihood behaves as $\sigma^{c(n-k)-k}$ as the scale
    $\sigma\to0$, for a shape-dependent exponent $c>0$.
  \item \emph{Spike}: the density concentrates entirely on one observation
    as the shape tends to an edge of its domain; no weakest bound exists,
    and the restriction simply ends the family at a named member.
  \item \emph{Branch}: the family changes character past a point (a
    discontinuous or non-monotone density on the far side).
  \item \emph{Endpoint}: an estimated endpoint approaches the sample
    extreme and the existence bound is attained there.
\end{itemize}

\subsection*{The collapse exponent: a weakest bound}

\begin{proposition}\label{prop:collapse}
Fix integers $n\ge3$ and $0\le k<n/2$, and suppose that as a scale
$\sigma\to0^+$ with a location held at an observation of multiplicity $k$,
$\ell(\sigma) = \{c(n-k)-k\}\log\sigma + O(1)$ for some $c>0$. Then:
(a) $\ell(\sigma)\to\infty$ as $\sigma\to0^+$ if and only if $c<k/(n-k)$;
(b) $\sup\{k/(n-k): n\ge3,\ 0\le k<n/2\} = 1$, not attained;
(c) $c\ge1$ excludes divergence for every such $(n,k)$, and it is the
weakest bound that does so for every $n$: for every $c<1$ there exist $n,k$
with $k<n/2$ for which $\ell(\sigma)\to\infty$.
\end{proposition}

\begin{proof}
(a) $\sigma^{c(n-k)-k}\to\infty$ as $\sigma\to0^+$ iff the exponent is
negative, i.e.\ $c(n-k)<k$, i.e.\ $c<k/(n-k)$ (using $n>k$).
(b) $x\mapsto x/(n-x)$ is strictly increasing on $[0,n)$, so for integer
$k<n/2$, $k/(n-k) < (n/2)/(n-n/2) = 1$. Along odd $n=2k+1$,
$k/(n-k)=k/(k+1)\to1$ as $k\to\infty$. So the supremum over all admissible
$(n,k)$ is $1$, not attained.
(c) If $c\ge1$, then $k/(n-k)<1\le c$ for every admissible pair, so by (a)
the path never diverges. If $c<1$, by (b) choose $k$ large and $n=2k+1$ so
that $k/(n-k)=k/(k+1)>c$; by (a) the path then diverges.
\end{proof}

Proposition~\ref{prop:collapse} is exactly the design's stated rule (R5):
``with tie multiplicity below $n/2$, $c\ge1$ is the weakest restriction
valid for every $n$.'' It applies directly where $c$ is itself one of the
free coefficients:

\begin{center}
\small
\begin{tabular}{@{}l l l l@{}}
\toprule
family & restricted coefficient & $c$ & exempt when held \\
\midrule
lst & \code{df} & $\mathrm{df}$ & \code{mu} \\
frechet & \code{lambda} & $\lambda$ & \code{mu} (inverse Weibull on $x-\mu$) \\
ht & \code{nu} & $\nu$ & \code{mu} \\
gev ($\xi>0$) & \code{xi} & $1/\xi$ & none (P0: diverges with \code{mu} or
  \code{sigma} held) \\
tlambda ($\lambda<0$) & \code{lambda} & $-1/\lambda$ & \code{mu} \\
hcauchy, slash & --- & $1$ (fixed) & none needed: $c=1$ already, so
  $k<n/2$ alone suffices \\
\bottomrule
\end{tabular}
\end{center}

Pareto II's free threshold, if it were supported, would need exactly
$c=\mathrm{shape}\ge1$ by the same proposition; it is the one family among
the newly added Pareto-type group whose collapse mechanism has no
independent endpoint singularity (below).

\subsection*{Spike, branch and endpoint}

\emph{Spike.} GED/SGED's density concentrates entirely on an observation as
shape $\nu\to0$. Profiling location at an observation of multiplicity $k$
and the scale, $\ell/n = -\log\{(n-k)/n\}/\nu + \tfrac12\log\nu + O(1) \to
\infty$ as $\nu\to0$ for \emph{every} nonconstant sample \cite{tumlinson2015}.
Unlike collapse there is no $n$-free weakest bound: the restriction
terminates the family at a named member ($\nu=1$, the Laplace; D6). P0's
amendment A3 sharpens the scope: with scale (\code{sd}) \emph{held}, the same
spike persists once the tied fraction exceeds $1-A/B = 0.1379185$,
$A=\tfrac32\log3$, $B=3\sqrt3/e$ (a Stirling expansion of the README
density), so the restriction must apply whenever \code{nu} itself is free,
regardless of what else is held, not only in the originally proposed
\code{while\_free: [mean, sd]} scope.

\emph{Branch.} Tukey's $\lambda$ family changes character at $\lambda=1$:
$Q'(u)=\sigma\{u^{\lambda-1}+(1-u)^{\lambda-1}\}$ gives a central-unimodal,
bounded-support density for $0<\lambda\le1$ and a discontinuous, U-shaped
density on $(1,2)$ whose boundary value at $\lambda=1$ is not the limit from
inside, i.e.\ not upper semicontinuous at that junction. The fitter
restricts to the principal branch $\lambda\le1$ (no exemption proved for any
hold) and accepts $\lambda=2$ only as an explicit alternate uniform alias
(scale differing by a factor of two from the $\lambda=1$ alias).

\emph{Endpoint.} For GEV/GPD with $\xi<-1$ the upper endpoint is finite
($\mu-\sigma/\xi$ for GEV); letting it approach the sample maximum from
above gives a $(-1-1/\xi)\log t$ term, $t\to0^+$, with negative coefficient,
hence $\ell\to\infty$ \cite{grimshaw1993}. At $\xi=-1$ exactly this
coefficient vanishes and the limit is the uniform (GPD) or
reversed-exponential (GEV) law; both \code{extraDistr::dgpd} and
\code{dgev} are \emph{zero} at that endpoint under their own (open-support)
convention, so the $\xi=-1$ supremum is literally unattained there.

\subsection*{Closure (R7a) and domains (R7)}

Rule R7a resolves the endpoint case: at a parameter-dependent support
endpoint the support is taken \emph{closed}, with the density there equal to
its limit from inside (the upper semicontinuous closure), for every family
whose reference defines an open support there (\code{dtnorm}, \code{dpower},
GPD/GEV at $\xi=-1$, the \code{tpois} lower bound). Estimated endpoints are
then genuinely attained and \code{dml} agrees with the fit; the earlier
hnorm-only special case becomes the general rule (implemented
\code{p3-closure}, 2026-10-28). With R7a in force, \code{p3-restrict}
implemented the shape bound at $\xi\ge-1$ for GPD and $\xi\in[-1,1]$ for
GEV exactly, superseding an intermediate design draft that had proposed
$\xi\ge-\tfrac12$.

Rule R7 separates two kinds of boundary case that the literal reference
density conflates: a continuous point mass (\code{sd=0}, beta shapes $0$ or
$\infty$) or an infinite parameter (\code{dht(nu=Inf)}, PIG \code{mean=Inf})
is excluded from the family's domain and stays an escape (``use
\code{mlX}''); a \emph{finite, discrete} point mass or other proper member
dominated by the family's base measure is included and silently attained
(NB \code{size=0}, \code{cmp} \code{lambda=0}/\code{nu=0}, \code{cmb}
\code{prob}$\in\{0,1\}$ and \code{size=0}, \code{bhatt} \code{a=0}, the
normal). \code{p3-domains} implemented this for NB, ZINB, PIG, CMP and CMB
(2026-10-29); bhatt's \code{a=0} normal member was already attained from P2.
\code{tpois}'s singleton-support boundary and a residual \code{zip}
\code{lambda=0} evaluation gap remain open (B7, \S\ref{sec:bugs}).

\subsection*{P0's counterexamples}

Each of P0's three design amendments (A3--A5) arose from a counterexample
the first-pass schema missed.

\emph{Independent endpoint singularities (A4).} Pareto III, Pareto IV and
Feller--Pareto each have a collapse mechanism at the threshold, with
exponent $c=\mathrm{shape}$ (Pareto III) or
$c=\mathrm{shape1}\cdot\mathrm{shape2}$ (Pareto IV, Feller--Pareto),
\emph{and} a second, independent endpoint singularity triggered whenever
$\mathrm{shape}<1$, $\mathrm{shape2}<1$, or
$\mathrm{shape2}\cdot\mathrm{shape3}<1$ respectively (Pareto II has no such
second singularity: its density is finite at the threshold for every
shape). A per-coefficient $c\ge1$ \code{fit\_domain} entry would stop the
collapse but not the endpoint singularity, so none of these three families,
nor \code{nsbeta} or \code{genbeta} (each singular below shape $1$ at an
end), supports an estimated threshold at all; \code{holds:} marks every
free-threshold subset \code{no-mle}.

\emph{The free GPD threshold} has no MLE for \emph{any} $\xi\in[-1,1]$:
with $\sigma$ held the likelihood strictly increases in $\mu$ up to
$\min(x)$ for $\xi>-1$, and at $\xi=-1$ it is a plateau that still excludes
the open interval $(\min,\max)$. A4 therefore leaves
\code{mlgpd(x, mu = NULL)} unsupported rather than proposing a shape bound.

\emph{Scope (A3).} GED's spike persists with \code{sd} held once tied mass
exceeds $13.79\%$ (above); GEV needs $\xi\le1$ even with \code{sigma} held,
and above the minimum with \code{mu} held; Fr\'echet needs $\lambda\ge1$
even with \code{sigma} held. The amendment: a restriction applies whenever
its own coefficient is free, independent of what else is held, with named
exemptions (\code{exempt\_when\_held}) only where P0 proves one (Fr\'echet's
\code{mu} held case is inverse Weibull on $x-\mu$ for every $\lambda>0$).

\emph{A well-posed restriction is not always a good default (A5).} Huber's
free-$\epsilon$ fit always exists mathematically (the family is
Hellinger-closed at both faces, normal and Laplace), yet the M4 simulation
(\S\ref{sec:new-families}) shows it answers ``use \code{mlnorm}'' for
$60$--$72\%$ of normal samples and up to $62\%$ of Cauchy samples at
$n=1000$. The fix is not a restriction but a kind change: \code{epsilon}
becomes \code{design}, required, not free by default.

\subsection*{P2b conditional skew fits}

The sstd and sged hold searches cover the previously unimplemented subsets.
Fixed-shape scale profiles and the sstd tie-collapse exponent are analytic;
the remaining joint and face comparisons use numerical searches, which do
not certify global optimality. Searches that exhaust their guards stop.
For sstd, holding \code{sd} alone does not prevent collapse: as free
$\nu\to2$, the ordinary $t$ scale shrinks even at fixed standard deviation.
Holding both \code{sd} and \code{nu} prevents this path. Holding the mean
is not holding the mode: their displacement depends on scale and shape,
including on limiting faces. The equality case of the sstd collapse criterion
and an sged uniform endpoint rate remain unresolved. These results and their
numerical qualifications are recorded in \code{dev/research/evaluations/2026-10-29-p2b-skew-report.md}.

\subsection*{P0's evidence}

Across the ten P0 jobs, $97$--$641$ finite package fits per group were
checked against an independent multistart oracle on the external reference
density (locscale 264/840, skewreal 227, positive 157, halfthreshold 108,
genbeta 641/1{,}961, extreme 212, bounded 97, counts 197/602, discrete2 337,
binomial 169 cases/total); every comparison agreed to $10^{-6}n$ except one
deliberate power-on-constant-data approximation. No group disputed the
parameter-kind table, and none of the bounded, discrete-count or
binomial-type groups needed a new R5 restriction, since every discrete
likelihood is $\le1$ and cannot diverge. Every job flags the same residual
limitation: multistart optimisation is evidence, not a certificate, of a
global maximum, and several exact criteria (GB2/genbeta joint faces,
sstd/sged global hold searches, beta-binomial free size) remain proved only
on subsets (\S\ref{sec:open}).

\section{Bugs found}\label{sec:bugs}

\begin{enumerate}
  \item \textbf{burr crash} (fixed, \code{411ea90}, \code{9290aac}). The
    log-moment start solves $\mu^2/\sigma^2 = r(k)$ on $k \in (0,1)$, where
    $r < 1$; mostly-Pareto data give a ratio $\ge 1$, the old clamp set
    $k = 10^{-15}$ and Newton drove $k$ to $0$, the pole of digamma.
  \item \textbf{Boost exceptions leak} (fixed, \code{0f02af1}). One
    no-throw policy at \code{special/special.hpp}; all Boost use goes
    through that file; a test checks nothing throws.
  \item \textbf{zip} returns $\lambda = 0$, $\pi = -\infty$ silently on
    ordinary zero-inflated samples (oracle $\hat\lambda = 1.42$,
    $\hat\pi = 0.47$). Fixed (\code{f1f1245}): the closed form
    $\hat\lambda = m + W_0(-m e^{-m})$ was evaluated on the $W_{-1}$ branch,
    whose value $-m$ gives the trivial root $\lambda = 0$.
  \item \textbf{cauchy} returns NaN for $n = 3$; \textbf{logis} returns NaN
    with one large outlier. Fixed (\code{776fc53}): safeguarded Newton,
    stable logistic likelihood; cauchy reports no MLE for $n < 3$ or when a
    value is repeated at least $n/2$ times.
  \item \textbf{binom} reports $\log L = -\infty$ at $\hat p = 1$ (the
    $0\log 0$ term); the correct value is $0$. Fixed (\code{f1f1245}),
    also in nbinom and geom.
  \item \textbf{std} reports a log-likelihood that disagrees with the
    reference density at very large $\nu$ (cancellation in the gamma terms).
    Open; it now matters less, since the normal face is detected.
  \item \textbf{Detection lives in the R adapter.} The nbinom
    underdispersion check existed only in the adapter, so a Python shell
    would silently lose it. Fixed: C++ status \code{Underdispersed}
    (\code{8cd0d92}).
  \item \textbf{sstd overflow} (fixed, \code{73f4fd5}). The standardising
    constant used $\Gamma(\tfrac{\nu-1}{2})/\Gamma(\tfrac{\nu}{2})$ with
    \code{tgamma}, which overflows for $\nu > 343$; the cost became NaN there
    and BFGS stopped at $\hat\nu = 343$ on normal data. Now \code{lgamma}.
  \item Fitters returned NaN instead of a clear error when an observation sits
    on a closed support bound (llogis, pareto, fatigue at $x = 0$). Fixed in
    the layout migration: those supports exclude $0$, so $x = 0$ is an input
    error.
\end{enumerate}

\subsection*{P0 bugs (TASK-79/85) and the TASK-77 scale-equivariance bug}

\begin{enumerate}
  \item[B1] \textbf{Holds silently ignored} (fixed, \code{3ca0b8d}). Partial
    argument matching let a misspelled or positional hold become a start
    value instead of a constraint: \code{mlweibull(x, shape = 2)} returned
    $\mathrm{shape}=2.08$ with nothing held. Also affected invweibull
    \code{shape}, lomax \code{lambda}, invpareto \code{scale}, kumar
    \code{a}, and the betapr/paralogis shapes. Fixed by requiring exact
    coefficient names.
  \item[B2] \textbf{Binomial \code{size} not the maximum} (fixed,
    \code{bugs85}). \code{mlbinom(c(2,6,6))} returned $\mathrm{size}=6$,
    $\ell=-6.827$, although size $10$ gives the unimodal maximum
    $\ell=-938.33$ on an earlier golden sample, and size $11$ beats $10$ by
    $0.055$ on an $n=100$ sample. Fixed with a bracketed adjacent-integer
    search using DeRiggi's unimodality theorem for the binomial profile
    likelihood \cite{deriggi1983}; the search still stops, rather than
    silently approximating, beyond size $10^6$ or on iteration exhaustion.
  \item[B3] \textbf{Boost exception escapes the no-throw policy} (fixed,
    \code{bugs85}). \code{mlkumar(rep(.4, 5))} crashed: a constant sample
    drove the L-moment start to a nonfinite tetragamma argument, and
    \code{boost::math::trunc} throws on NaN even though digamma/trigamma
    return NaN quietly. Fixed by detecting constant samples before the
    L-moment step and guarding nonfinite Newton inputs.
  \item[B4] \textbf{$\ell=-\infty$ returned instead of a diagnostic} (fixed,
    \code{bugs85}). \code{mlrayleigh(c(0, 0, 1, 2))}: an observation at zero
    has density zero for every scale, so the correct response is a status,
    not a numeric fit.
  \item[B5] \textbf{Attained MLE reported as non-convergence} (fixed,
    \code{bugs85}). \code{mltriang(rep(.4, 5), c = -1)} (also with
    \code{b = 2} held): the closed-form optimum $a=c=-1$, $b=1.8$ (or $2$)
    was correct, but the endpoint solver reported iteration failure at
    every tested tolerance and iteration limit; fixed by preserving
    equality between a free endpoint and a held mode when undoing
    standardization.
  \item[B6] \textbf{Tie rule stricter than needed} (fixed, \code{bugs85}).
    GEV and Fr\'echet counted the \emph{largest} tie instead of the tie at
    the minimum when deciding whether the fixed-shape collapse applies,
    rejecting some samples whose MLE exists; fixed to count the minimum's
    own multiplicity and reject exactly when $\xi k \ge n-k$. The equality
    case $k=n/2$ now compares the interior with the collapse face.
  \item[B7] \textbf{Attained boundary maxima rejected} (partly fixed,
    \code{p3-domains}). cmp's point-mass ($\lambda=0$) and geometric
    ($\nu=0$, $\lambda<1$) faces, NB/ZINB \code{size=0}, PIG
    \code{shape=0} and cmb's endpoint/size-0 members are now returned as
    attained R7 members rather than escapes (\S\ref{sec:holds-restrictions}).
    tpois's singleton-support boundary is still rejected.
  \item[B8] \textbf{README domains narrower than the reference} (fixed,
    \code{p3-domains}). \code{dcmb} accepts $\mathrm{prob}\in\{0,1\}$ and
    size up to $10^7$; \code{dbhatt} accepts $a=0$ (the normal, already
    attained from P2); the domain audit widened NB, ZINB, PIG, CMP and CMB
    to match.
  \item[TASK-77] \textbf{\code{std} returned its own start values on
    daily-return data} (fixed, \code{scaleq}). On DAX log returns (sample
    sd $\approx0.01$), the raw-unit BFGS gradient had norm $\approx2610$;
    neither of its two Armijo line searches accepted a step, giving
    \code{LineSearchFailure} at iteration 1 with $\nu$ and $\ell$
    \emph{unchanged} from the start values: no convergence rule fired, and
    the adapter's old policy (return whatever BFGS landed on) silently
    accepted that failure as a fit. The old normalization guard (skip
    centering below sample sd $10^{-3}$) missed this data. Fixed by always
    centering and dividing by the sample sd before every interior/face
    comparison, restoring the Jacobian afterward; seven other fitters
    (sstd, gompertz, power, fatigue, paralogis/invparalogis, betapr, lomax)
    had the same unit dependence and were fixed the same way. This also
    replaced the nlminb-legacy wrapper rule with R6
    (\S\ref{sec:holds-restrictions}): a failed iterative optimizer now
    stops unless its returned point passes a score gate, never silently
    returning an unchecked last iterate.
\end{enumerate}

\section{Open questions and decisions}\label{sec:open}

\begin{enumerate}
  \item D1, escape policy: decided, stop with ``use \code{mlX}''; every
    detected escape is declared under \code{limits:} in the distribution's
    README.
  \item D5, nakagami: decided, allow $m > 0$.
  \item D6, ged and sged: decided and implemented, $\nu \ge 1$. (fGarch's
    own \code{gedFit} uses the ill-posed lower bound $0$.)
  \item D7, limit models: decided and implemented, hnorm and lst; messages
    only for half-$t$, two-piece $t_2$ and half-GED.
  \item snorm, sstd: done. sstd's comparison costs about 20\,ms at
    $n = 500$ (seven interior starts, two half-$t$ searches, sixteen $t_2$
    starts); it is a multistart comparison, not a certificate.
  \item Fixed parameters (\code{dev/roadmap/2026-05-08-variable-parameter-selection.md}):
    \code{scale} in burr, invburr, betapr, paralogis; \code{nu} in std,
    \code{df} in lst; \code{size} in binom, nbinom; fatigue's $\mu$ always.
    Limits that need a held parameter to move carry \code{fixed:} in the
    README.
  \item Open: a half-$t$ model (would give sstd's half-$t$ face a pointer).
  \item R1--R8, parameter rules (\code{dev/roadmap/2026-10-29-parameter-rules.md}):
    decided 2026-10-29 and implemented through P3; \S\ref{sec:holds-restrictions}
    gives the mechanisms. D1 and D7 above are now instances of R6 and the
    \code{limits:} schema generally.
  \item GEV/GPD shape bound: an intermediate design draft proposed
    $\xi\ge-\tfrac12$; P3 implements the originally intended $\xi\ge-1$
    (GPD) and $\xi\in[-1,1]$ (GEV), attained via R7a closure.
\end{enumerate}

\subsection*{Open mathematics}

\begin{itemize}
  \item Certified global search: every new multistart/grid fitter
    (half-$t$ df, discretized-gamma shape, Skellam, slash, Bhattacharjee,
    Tukey lambda, generalized beta, GB2, \ldots) is evidence, not a
    certificate, of a global maximum.
  \item dgamma: is the discretized-gamma shape profile unimodal?
  \item Skellam: uniqueness of the joint maximum.
  \item Slash: a closed sample condition for the exactly-half-tied case
    (the implemented criterion is variational, comparing the interior
    profile to $B_v$).
  \item GB2 and generalized beta: is the catalogue of limit models
    complete, and what is the exact classification of the remaining hold
    subsets?
  \item Tukey lambda: the U-shaped branch $1<\lambda<2$ is not fitted; the
    uniform junction at $\lambda=1$ is not upper semicontinuous.
  \item Negative hypergeometric: the inverse-sampling design (unknown
    population size) is unsupported.
  \item Truncated binomial: free-size searches are implemented, but their
    global integer-profile search is numerical and not certified.
  \item Free-threshold fits with another parameter held (Pareto II--IV,
    Feller--Pareto, nonstandard beta): each subset is its own open
    existence question (A4).
  \item Translated half-$t$ uses the collapse restriction $\nu\ge1$
    when its threshold is free.
  \item Beta-binomial with fixed size: existence proved for size $2$, open
    for size $\ge3$.
  \item sstd and sged: a proof that the listed
    translated half-family faces exhaust every joint escape.
  \item Beta-negative-binomial: an explicit sample statistic for the
    NB/Poisson face criterion (only the implicit ``some finite point beats
    every face'' criterion is proved).
  \item Fatigue: the $3k=n$ endpoint profile comparison.
  \item Beta prime: free-scale closure completeness.
  \item GEV and GPD on $[-1,-\tfrac12]$: interior, endpoint and
    upper-grid profile maxima are all observed numerically; roots are not
    certified.
  \item sn: the escape-frequency discrepancy with
    \cite{azzalini1985,sartori2006} at $\alpha=5,10$, $n=50$
    (\S\ref{sec:new-families}) is unresolved.
\end{itemize}

\begin{thebibliography}{99}

\bibitem{fs1998} Fern\'andez, C. and Steel, M. F. J. (1998). On Bayesian
modeling of fat tails and skewness. \emph{Journal of the American
Statistical Association}, 93(441), 359--371.

\bibitem{kk2003} Kleiber, C. and Kotz, S. (2003). \emph{Statistical Size
Distributions in Economics and Actuarial Sciences}. Wiley.

\bibitem{tumlinson2015} Tumlinson, S. E. (2015). On the non-existence of
maximum likelihood estimates for the extended exponential power distribution
and its generalizations. \emph{Statistics \& Probability Letters}, 107,
111--114. \href{https://doi.org/10.1016/j.spl.2015.08.009}{doi:10.1016/j.spl.2015.08.009}.

\bibitem{shao2004} Shao, Q. (2004). Notes on maximum likelihood estimation
for the three-parameter Burr XII distribution. \emph{Computational Statistics
\& Data Analysis}, 45(3), 675--687.
\href{https://doi.org/10.1016/S0167-9473(02)00367-5}{doi:10.1016/S0167-9473(02)00367-5}.

\bibitem{fs1999} Fern\'andez, C. and Steel, M. F. J. (1999). Multivariate
Student-$t$ regression models: pitfalls and inference. \emph{Biometrika},
86(1), 153--167. \href{https://doi.org/10.1093/biomet/86.1.153}{doi:10.1093/biomet/86.1.153}.

\bibitem{delcastillo2009} del Castillo, J. and Daoudi, J. (2009). Estimation
of the generalized Pareto distribution. \emph{Statistics \& Probability
Letters}, 79(5), 684--688.
\href{https://doi.org/10.1016/j.spl.2008.10.021}{doi:10.1016/j.spl.2008.10.021}.

\bibitem{grimshaw1993} Grimshaw, S. D. (1993). Computing maximum likelihood
estimates for the generalized Pareto distribution. \emph{Technometrics},
35(2), 185--191.
\href{https://doi.org/10.1080/00401706.1993.10485040}{doi:10.1080/00401706.1993.10485040}.

\bibitem{corless1996} Corless, R. M., Gonnet, G. H., Hare, D. E. G., Jeffrey,
D. J. and Knuth, D. E. (1996). On the Lambert $W$ function. \emph{Advances in
Computational Mathematics}, 5(1), 329--359.
\href{https://doi.org/10.1007/BF02124750}{doi:10.1007/BF02124750}.

\bibitem{renext} Deville, Y. (2025). \emph{Renext: Renewal Method
for Extreme Values Extrapolation}. R package version 3.1-5.
\href{https://doi.org/10.32614/CRAN.package.Renext}{doi:10.32614/CRAN.package.Renext}.

\bibitem{deriggi1983} DeRiggi, D. F. (1983). Unimodality of likelihood
functions for the binomial distribution. \emph{Journal of the American
Statistical Association}, 78(381), 181--183.
\href{https://doi.org/10.1080/01621459.1983.10477949}{doi:10.1080/01621459.1983.10477949}.

\bibitem{burridge1981} Burridge, J. (1981). A note on maximum likelihood
estimation for regression models using grouped data. \emph{Journal of the
Royal Statistical Society, Series B}, 43(1), 41--45.
\href{https://doi.org/10.1111/j.2517-6161.1981.tb01146.x}{doi:10.1111/j.2517-6161.1981.tb01146.x}.

\bibitem{azzalini1985} Azzalini, A. (1985). A class of distributions which
includes the normal ones. \emph{Scandinavian Journal of Statistics}, 12,
171--178.

\bibitem{sartori2006} Sartori, N. (2006). Bias prevention of maximum
likelihood estimates for scalar skew normal and skew $t$ distributions.
\emph{Journal of Statistical Planning and Inference}, 136, 4259--4275.

\bibitem{kadane2016} Kadane, J. B. (2016). Sums of possibly associated
Bernoulli variables: the Conway--Maxwell--binomial distribution.
\emph{Bayesian Analysis}, 11(2), 403--420.
\href{https://doi.org/10.1214/15-BA955}{doi:10.1214/15-BA955}.

\bibitem{irwin1975} Irwin, J. O. (1975). The generalized Waring
distribution, Part I. \emph{Journal of the Royal Statistical Society,
Series A}, 138(1), 18--31.
\href{https://doi.org/10.2307/2345247}{doi:10.2307/2345247}.

\bibitem{xekalaki1984} Xekalaki, E. (1984). The bivariate generalized
Waring distribution and its application to accident theory. \emph{Journal
of the Royal Statistical Society, Series A}, 147(3), 488--498.
\href{https://doi.org/10.2307/2981580}{doi:10.2307/2981580}.

\bibitem{berrymielke1998} Berry, K. J. and Mielke, P. W. (1998). The
negative hypergeometric probability distribution: sampling without
replacement from a finite population. \emph{Perceptual and Motor Skills},
86(1), 207--210.

\end{thebibliography}

\end{document}
